%%%%%%%%%%  scan idx 664  |  folio 651  %%%%%%%%%%
\SGAPrintPage{651}%
still integral over \(S\), with fibers of dimension \(2\); by 4.3, we may replace \(X\) by \(Y\). Finally, by 4.3 and 3.10 (using resolution of singularities for surfaces), after replacing \(S\) by a flat covering of a nonempty open of \(S\), we may replace \(X\) by a smooth \(S\)-scheme.

Now, by 4.6 (i) \(\Longleftrightarrow\) (iii) and the finiteness and specialization theorem for \'etale cohomology groups (SGA 4 XVI 2.2), 5.1 follows from the next proposition.

\sgaStableTarget{sga6.e13.prop.5-2}\textbf{Proposition 5.2.} \emph{Let \(k\) be an algebraically closed field and \(X\) a proper smooth \(k\)-scheme. Then the group \(C^p_{\rm num}(X)\) of algebraic cycles of codimension \(p\) on \(X\), modulo numerical equivalence, is free of rank \(\le b^{2p}\), the Betti number.}

Indeed, if we forget the action of the roots of 1, \'etale cohomology \(H^*(X,\Q_\ell)=_{\rm dfn}\left[\varprojlim H^*(X,\Z/\ell^\nu\Z)\right]\otimes_{\Z_\ell}\Q_\ell\) where \((\ell,\operatorname{car}(k))=1\) becomes a Weil cohomology: it satisfies Poincar\'e duality and the K\"unneth formula; and it receives a functorial, nonzero homomorphism from the ring of algebraic cycles modulo rational equivalence \(C^*_{\rm rat}(X)\) (see X 7.13).

Let \(a_1,\ldots,a_m\in C^{r-p}_{\rm rat}(X)\), where \(r=\dim(X)\), be elements that map to a basis of the \(\Q_\ell\)-vector subspace of \(H^{2(r-p)}(X,\Q_\ell)\) generated by algebraic cycles. Consider the homomorphism
\[
   \alpha:C^p_{\rm rat}(X)\longrightarrow \Z^m
\]
defined by \(\alpha(b)=(\langle b.a_1\rangle,\ldots,\langle b,a_m\rangle)\), where \(\langle\ \rangle\) denotes the intersection number. Then one readily sees that the image of \(\alpha\) is \(C^p_{\rm num}(X)\).

%%%%%%%%%%  scan idx 665  |  folio 652  %%%%%%%%%%
\SGAPrintPage{652}%
\sgaStableTarget{sga6.e13.rem.5-3}\textbf{Remark 5.3.} Let \(k\) be an algebraically closed field, and \(X\) a proper smooth surface over \(k\). Then the second Betti number \(b^2(X)\) can be interpreted as \(c_2(X)-2+2b^1(X)\), where \(b^1(X)=2\dim(\Pic_{X/k})\), and \(c_2(X)\) = the second Chern class = the Euler-Poincar\'e characteristic = the intersection number of the diagonal of \(X\times X\) with itself; thus \(b^2(X)\) has an interpretation independent of cohomology. The inequality \(\rang(C^1_{\rm num}(X))\le b^2(X)\) in this form for \(\operatorname{car}(k)>0\) is due (before \'etale cohomology) to IGUSA, who used deep results on the structure of the tame fundamental group of an algebraic curve minus a few points.

\sgaStableTarget{sga6.e13.thm.5-4}\textbf{Theorem 5.4.} \emph{Let \(k\) be an algebraically closed field, and \(X\) a proper \(k\)-scheme. Then there exist finitely many closed integral curves \(Y_1,\ldots,Y_p\) in \(X\) with the following property: a subset \(\Lambda\) of the scheme \(\Pic_{X/k}\) is quasi-compact if and only if the intersection numbers \(\langle\lambda.Y_i\rangle\) for \(\lambda\in\Lambda\) remain bounded. Moreover, we may take for \(p\) the rank of the N\'eron-Severi group \(\Pic_{X/k}(k)/\Pic^0_{X/k}(k)\).}

Indeed, since \(\Pic^\tau_{X/k}\) is an open quasi-compact subgroup of \(\Pic_{X/k}\) (4.7), \(\Lambda\) is quasi-compact if and only if the corresponding set of points in the vector space \(\bigl[\Pic_{X/k}(k)/\Pic^\tau_{X/k}(k)\bigr]\otimes\Q\) is finite. But \(\tau\)-equivalence is equal to numerical equivalence (4.6 (i) \(\Longleftrightarrow\) (iii)), so we may choose among the curves in \(X\) curves that map to a basis of the dual space; hence the assertion.

%%%%%%%%%%  scan idx 666  |  folio 653  %%%%%%%%%%
\SGAPrintPage{653}%
\sgaStableTarget{sga6.e13.sec.6}\subsection*{6. Appendix: Study of $(b)$-sheaves on $P=\mathbb P^N_k$.}

\sgaStableTarget{sga6.e13.para.6-1}\textbf{6.1.} Let $k$ be an algebraically closed field, and let $P=\mathbb P^N_k$. Let $F$ be a coherent sheaf on $P$, with Hilbert polynomial
\[
   \chi(F(n))=\sum_{i=0}^{r} a_i \binom{n+i}{i},
\]
and let
\[
   (c)=(c_0,\ldots,c_r)
\]
be a sequence of integers $\ge (a)$. Put
\[
   a_{m,i}=\sum_{j=0}^{r-i}a_{j+i}\binom{m-1+j}{j}\quad\text{and}\quad c_{m,i}=\sum_{j=0}^{r-i}c_{j+i}\binom{m-1+j}{j}.
\]
Finally, for every integer $q\ge1$, denote the largest coherent subsheaf of $F$ with support of dimension $\le q-1$ by $N_q$, and put $F_q=F/N_q$.

\sgaStableTarget{sga6.e13.lem.6-2}\textbf{Lemma 6.2.} \emph{Under the conditions of 6.1, consider an exact sequence}
\[
   0\longrightarrow I\longrightarrow {O}_P^{\oplus M}\longrightarrow F(m)\longrightarrow0
\]
\emph{with $m\ge0$. Then $I$ is $p$-regular with \(p=P_r(c_{m,0},\ldots,c_{m,r})\), where $P_r$ is the $r$-th $(b)$-polynomial (1.9).}

Indeed, ${O}_P$ is clearly a $(0,\ldots,0,1)$-sheaf ($N+1$ terms), so $I$ is a $(0,\ldots,0,M)$-sheaf by 1.6. By 2.10, \(\chi(I(n))=M\binom{n+N}{N}-\sum_{i=0}^{r} a_{m,i}\binom{n+i}{i}\). Thus the assertion follows from 1.11 and 1.10.

\sgaStableTarget{sga6.e13.prop.6-3}\textbf{Proposition 6.3.} \emph{Under the conditions of 6.1, suppose that $F$ is a $(b)$-sheaf. Then one can calculate, by means of universal polynomials in the $a_i$, $b_i$ and $N$, integers $m, M, m_1$ and $M_1$ such that there is an exact sequence}
\[
   {O}_P(-m_1)^{\oplus M_1}\longrightarrow
   {O}_P(-m)^{\oplus M}\longrightarrow F\longrightarrow0.
\]
\emph{Moreover, $N$ occurs only in the calculation of $M_1$.}

%%%%%%%%%%  scan idx 667  |  folio 654  %%%%%%%%%%
\SGAPrintPage{654}%
Indeed, take $m=P_{r-1}(c_0,\ldots,c_{r-1})$ with $c_i=b_i-a_i$, and $M=\sum a_i\binom{m+i}{i}$, so that there is an exact sequence
\[
   0\longrightarrow I\longrightarrow {O}_P^{\oplus M}
   \longrightarrow F(m)\longrightarrow0
\]
by 1.11 and 1.3. Similarly, taking 6.2 into account, take $m_1=m+p$ with $p=P_r(a_{m,0},\ldots,a_{m,r})$, and $M_1=M\cdot\binom{p+N}{N}-\sum a_i\binom{m_1+i}{i}\quad(=\chi(I(p)))$.

\sgaStableTarget{sga6.e13.thm.6-4}\textbf{Theorem 6.4.} \emph{Under the conditions of 6.1, suppose that $F$ is a quotient of ${O}_P(-m)^{\oplus M}$ with $m\ge0$. Let $b_i=P_{r-i}(c_{m,i},\ldots,c_{m,r})$ for $i=0,\ldots,r$, where $P_s$ is the $s$-th $(b)$-polynomial (1.9); let $b=b_1+m-1$, and let $B=\sum_{i=1}^{r} a_i\binom{b+i}{i}$. Then}:
\begin{enumerate}[label=(\roman*)]
\item \emph{$F$ is $b$-regular.}
\item \emph{$-B\le a_0=h^0(F(b))-B$}
\item \emph{$F$ is a $(b_0,\ldots,b_r)$-sheaf.}
\end{enumerate}

Indeed, take a section $\sigma$ of ${O}_P(1)$ such that we obtain a commutative diagram with exact rows and columns
\[
\begin{tikzcd}[column sep=large,row sep=large]
&0 \arrow[d]&0 \arrow[d]&0 \arrow[d]&\\
0 \arrow[r]&I(-1) \arrow[r] \arrow[d]&I \arrow[r] \arrow[d]&J \arrow[r] \arrow[d]&0\\
0 \arrow[r]&{O}_P(-1)^{\oplus M} \arrow[r] \arrow[d]&{O}_P^{\oplus M} \arrow[r] \arrow[d]&{O}_Y^{\oplus M} \arrow[r] \arrow[d]&0\\
0 \arrow[r]&F(m-1) \arrow[r] \arrow[d]&F(m) \arrow[r] \arrow[d]&G(m) \arrow[r] \arrow[d]&0\\
&0&0&0&.
\end{tikzcd}
\]
%%%%%%%%%%  scan idx 668  |  folio 655  %%%%%%%%%%
\SGAPrintPage{655}%
Since \(\chi(G(n))=\sum_{i=0}^{r-1} a_{i+1}\binom{n+i}{i}\) by 1.7, $J$ is $b_1$-regular by 6.2; thus \(H^q(I(p))=0\) for $q\ge2$ and $p\ge b_1-q$ by 1.4 (i). The exact sequence \(H^q({O}_P^{\oplus M}(p))\longrightarrow H^q(F(m+p))\longrightarrow H^{q+1}(I(p))\) then shows that $F(m)$ is $(b_1-1)$-regular; hence (i).

Assertion (i) implies that \(0\le h^0(F(b))=\chi(F(b))=\sum_{i=0}^{r}a_i\binom{b+i}{i}\); hence (ii). By 2.10 and 1.9,
\[
   \chi(F(b))=\sum_{i=0}^{r} a_{m,i}\binom{b_1-1+i}{i}\le b_1+\sum_{i=0}^{r} c_{m,i}\binom{b_1-1+i}{i}=b_0\,;
\]
thus, $h^0(F(-1))\le h^0(F(b))\le b_0$. By induction on $r$, we may suppose that $G$ is a $(b_1,\ldots,b_r)$-sheaf; hence (iii) by 1.6.2.

\sgaStableTarget{sga6.e13.lem.6-5}\textbf{Lemma 6.5.} \emph{Under the conditions 6.1, suppose that $F$ has no associated prime cycle of dimension $0$.}
\begin{enumerate}[label=(\roman*)]
\item \emph{If $h^0(F)\ge1$, then $h^0(F(-1))\le h^0(F)-1$.}
\item \emph{$H^0(F(-n))=0$ for $n\ge h^0(F)$.}
\item \emph{Let an exact sequence $0\to F(-1)\to F\to G\to0$ and an integer $n_0\ge0$ satisfy $H^0(G(-n_0))=0$; then $H^0(F(-n_0))=0$ and $h^0(F)\le n_0.h^0(G)$.}
\end{enumerate}
Indeed, let $\sigma\in H^0(F)$, $\sigma\ne0$; then $\sigma$ defines a nonzero subsheaf $F'={O}_P.\sigma$ of $F$, so $s=\dim.\operatorname{supp}(F')$ is $\ge1$. Take a commutative diagram with exact rows and columns
\[
\begin{tikzcd}[column sep=large,row sep=large]
&0 \arrow[d]&0 \arrow[d]&0 \arrow[d]&\\
0 \arrow[r]&F'(-1) \arrow[r] \arrow[d]&F' \arrow[r] \arrow[d]&G' \arrow[r] \arrow[d]&0\\
0 \arrow[r]&F(-1) \arrow[r]&F \arrow[r]&G \arrow[r]&0 & .
\end{tikzcd}
\]
